\documentclass{article}
\usepackage{graphicx}
\usepackage{amsthm,amsmath,amssymb,amsfonts,amscd}
\usepackage[utf8x]{inputenc}
\usepackage{tikz-cd} 
\usepackage{float}
\usepackage{adjustbox}
\usepackage{tikz-3dplot}
\usepackage[backend=biber,style=alphabetic]{biblatex}
\usepackage{xcolor}
\addbibresource{bibliography.bib}
\newtheorem*{theoremA}{Theorem A}
\newtheorem{theorem}{Theorem}[section]

\newtheorem{proposition}[theorem]{Proposition}
\newtheorem{lemma}[theorem]{Lemma}
\newtheorem{corollary}[theorem]{Corollary}
\newtheorem{conjecture}[theorem]{Conjecture}
\newtheorem{problem}[theorem]{Problem}
\newtheorem{algorithm}[theorem]{Algorithm}

\theoremstyle{definition}
\newtheorem{definition}[theorem]{Definition}
\newtheorem{example}[theorem]{Example}

\theoremstyle{remark}
\newtheorem{remark}[theorem]{Remark}
\newtheorem*{remarknonumber}{Remark}
\newtheorem{observation}[theorem]{Observation}

\theoremstyle{theorem}
\newtheorem*{newmanslemma}{Newman's Lemma}

\theoremstyle{definition}
\newtheorem{assumption}[theorem]{Assumption}

\title{The Word problem on Partially 2-nilpotent
groups}
\author{Javier Mejia}
\date{Feb 2026}


\begin{document}
    
\maketitle

\begin{section}{Definition and Basic Properties}
    Previously we defined the partially $\omega$ groups and presented the example of the partially 2-nilpotent groups, when $\omega = [[x,y],z]$.

\begin{definition}
    Let $V$ be a set and $\chi=(V,E\subseteq P(V))$ be a hypergraph on $V$. Define the group

   $$G(\chi)=\langle V \mid [[g_1,g_2],g_3]=1\;\forall g_1,g_2,g_3\in \langle e\rangle \;\forall e\in E\subseteq P(V) \rangle$$

    called the \textbf{partially 2-nilpotent group} of $\chi$.
\end{definition}

The generators are the vertices of the hypergraph, and the relations state
that elements generated by any hyperedge have a 2-nilpotent relation.

\begin{assumption}\label{assump:dense}
Throughout this paper, we assume all hypergraphs $\chi = (V, E)$ are \textbf{dense} (also called \textbf{downward-closed}), meaning:
$$\text{If } e \in E \text{ and } e' \subseteq e , \text{ then } e' \in E$$

In other words, every subset of a hyperedge is itself a hyperedge.
\end{assumption}

\begin{remark}
To work effectively with these groups, we introduce additional structure. Let $V = \{v_1, \ldots, v_n\}$ with a total order $v_1 < v_2 < \cdots < v_n$. We decompose the hyperedge set as:
$$E = E_1\sqcup E_2 \sqcup E_3\sqcup \cdots \sqcup E_{|V|}$$ 
where $E_k=\{\text{hyperedges of }\chi \text{ with exactly }k \text{ vertices} \}$.
\end{remark}

For partially 2-nilpotent groups, $E_1$ (singleton edges) produces only trivial relations $[[x,x],x]=1$, so we omit them without loss of generality.

\subsection{Extended Hyperedges and the Augmented Alphabet}

\begin{definition}
For a hyperedge $E = \{v_{i_1}, \ldots, v_{i_r}\} \in \mathcal{E}$ (written in order $i_1 < \cdots < i_r$), define the \textbf{extended hyperedge}:
$$\bar{E} = \{v_{i_1}, \ldots, v_{i_r}\} \cup \{e_{i_j i_k} \mid 1 \leq j < k \leq r\}$$
where $e_{ij} = [v_i, v_j]$ (with $i < j$) is a new symbol that denotes the commutator generator.

The \textbf{augmented alphabet} is:
$$A = V \cup E_2 = V \cup \{e_{ij} \mid \{v_i, v_j\} \in E_2, i<j\}$$
\end{definition}

\begin{remark}
The density assumption ensures that if $\{v_1, v_2, v_3\} \in E_3$, then $\{v_1, v_2\}, \{v_1, v_3\}, \{v_2, v_3\} \in E_2$. This guarantees that all pairwise commutators $[v_i, v_j]$ for vertices in a common hyperedge become generators in $A = V \cup E_2$ (see Definition~1.5 above).
\end{remark}

\begin{definition}
For the alphabet $A = V \cup E_2$, define:
$$A^{\pm} = \{a^{+1}, a^{-1} \mid a \in A\}$$
the set of formal symbols consisting of elements of $A$ with exponents $\pm 1$. 

We adopt the notational conventions:
\begin{itemize}
    \item $a^{+1}$ is written as $a$
    \item $(a^{-1})^{-1} = a$
    \item For $x = a^\epsilon \in A^{\pm}$, we write $x^{-1} = a^{-\epsilon}$
    \item In the total order on $A^{\pm}$, we set $a^{-1}<a $ for all $a \in A$
\end{itemize}

A word in $A^{\pm}$ is an element of the free monoid $(A^{\pm})^*$.
\end{definition}

The augmented alphabet $A$ extends the vertex set to include edge generators representing commutators. We extend the total order on $V$ to $A$ by setting:
$$v_1 < v_2 < \cdots < v_n < e_{11} < e_{12} < \cdots < e_{1n} < e_{23} < \cdots < e_{nn}$$
The full total order on $A^{\pm}$ is then defined by $a^{-1} < a$ for all $a \in A$, and $a^{-1} < b^{-1}$ if and only if $a < b$. Concretely:
$$v_1^{-1} < v_1 < v_2^{-1} < v_2 < \cdots < v_n^{-1} < v_n < e_{11}^{-1} < e_{11} < e_{12}^{-1} < e_{12} < \cdots < e_{nn}^{-1} < e_{nn}$$

\begin{remark}
The extended hyperedge $\bar{E}$ tells us which elements can ``interact'' --- that is, which generators can be moved past each other (possibly generating commutators) within the context of hyperedge $E$.
\end{remark}

\begin{lemma}[Property of 2-nilpotent groups]\label{lem:2nilpotent}
    Let $G$ be a group generated by a set $S$. If $[[s_i, s_j], s_k] = 1$
    for all $s_i, s_j, s_k \in S$, then $G$ is 2-nilpotent: $[[g_1,g_2],g_3]=1$
    for all $g_1,g_2,g_3 \in G$.

    In particular, if $x,y \in G$ satisfy $[[x,y],x]=1$ and $[[x,y],y]=1$,
    then $\langle x,y\rangle$ is 2-nilpotent.
\end{lemma}

\begin{proof}
    We show directly that every commutator generator $[s_i, s_j]$ (for
    $s_i, s_j \in S$) is central in $G$, i.e., $[[s_i, s_j], g] = 1$ for
    all $g \in G$. We induct on the word length $|g|$ in $S^{\pm 1}$.

    \textit{Base case ($|g|=1$):} If $g = s_k \in S$, then
    $[[s_i,s_j],s_k]=1$ by hypothesis. If $g = s_k^{-1}$, then
   $$[[s_i,s_j],\,s_k^{-1}] \;=\; s_k\,[[s_i,s_j],s_k]^{-1}\,s_k^{-1}
      \;=\; s_k \cdot 1 \cdot s_k^{-1} \;=\; 1.$$

    \textit{Inductive step:} Write $g = h \cdot t$ with $t \in S^{\pm 1}$
    and $|h| < |g|$. Apply the exact identity (valid in any group \cite{MKS})
    \begin{equation}\label{eq:exact_commutator}
      [x,\,ab] \;=\; [x,b]\cdot[x,a]^b,
    \end{equation}
    with $x=[s_i,s_j]$, $a=h$, $b=t$:
   $$[[s_i,s_j],\,h\cdot t] \;=\; [[s_i,s_j],\,t]\cdot[[s_i,s_j],\,h]^t.$$
    The first factor equals $1$ by the base case ($|t|=1$). The second
    equals $1^t = 1$ by the inductive hypothesis on $|h| < |g|$.

    Since every generator $[s_i,s_j]$ of $\gamma_2(G)$ is central in $G$,
    we have $\gamma_2(G)\subseteq Z(G)$. Therefore
    $\gamma_3(G)=[\gamma_2(G),G]=1$, so $G$ is 2-nilpotent.

    For the ``in particular'' clause: with $S=\{x,y\}$, the hypotheses
    $[[x,y],x]=1$ and $[[x,y],y]=1$ imply all weight-3 generator
    commutators vanish---the remaining cases $[[y,x],x]$ and $[[y,x],y]$
    follow since $[y,x]=[x,y]^{-1}$ and the hypothesis gives
    $[[x,y],x]=[[x,y],y]=1$---so the main result applies.
\end{proof}

\subsection{Alternative Presentation}

\begin{lemma}\label{lem:alternative}
    The following alternative definition of partially 2-nilpotent groups is equivalent to Definition 1.1:

    Let $E = E_1\sqcup E_2\sqcup E_3\sqcup \cdots \sqcup E_{|V|}$ be the hyperedge decomposition of $\chi$. Define 
    \begin{align*}
    G'(\chi)=\langle V\sqcup E_2 \mid
    &[[x,y],z]=1 \quad \forall x,y,z \in e, \; e \in E,\\
    &[x,y]=e_{xy} \quad \forall \{x,y\} \in E_2 \rangle
    \end{align*}
    
    Then $G'(\chi) \cong G(\chi)$.
\end{lemma}

\begin{proof}
We show the presentations are equivalent using Tietze transformations.

Starting with:
$$G(\chi)=\langle V \mid [[g_1,g_2],g_3]=1\;\forall g_1,g_2,g_3\in \langle e\rangle, \; \forall e\in E \rangle$$

\textit{Step 1 (Tietze I):} For each $\{x,y\} \in E_2$ with $x < y$ in the fixed total order on $V$, simultaneously introduce the new generator $e_{xy}$ together with its defining relation $[x,y] = e_{xy}$. Since $e_{xy}$ is defined as a word in the existing generators, this is a redundant generator addition; by Tietze I the resulting presentation is equivalent to $G(\chi)$. This gives:
$$G_1 = \langle V \sqcup E_2 \mid [[g_1,g_2],g_3]=1\;\forall g_1,g_2,g_3\in \langle e\rangle,\; \forall e\in E,\; [x,y]=e_{xy}\;\forall\{x,y\}\in E_2 \rangle$$

\textit{Step 2 (Replace general relations with specialized versions):} 

We claim the general relation $[[g_1,g_2],g_3]=1$ for all $g_1,g_2,g_3\in \langle e\rangle$ is equivalent to:
$[[x,y],z]=1$ for all $x,y,z \in e$ where $e \in E$

\textit{($\Rightarrow$)} Immediate by specializing $g_1, g_2, g_3$.

\textit{($\Leftarrow$)} We must show $\langle e \rangle$ is 2-nilpotent for any $e \in E$. The generators of $\langle e\rangle$ are the vertices of $e$, and by relations (1) and (2), the weight-3 commutator $[[s_i,s_j],s_k]$ vanishes for all $s_i,s_j,s_k \in e$. By Lemma~\ref{lem:2nilpotent} (applied with $S = e$), $\langle e \rangle$ is 2-nilpotent, so $[[g_1,g_2],g_3]=1$ for all $g_1,g_2,g_3 \in \langle e \rangle$.

Thus $G_1 = G'(\chi)$.
\end{proof}

\begin{example}
A two-vertex hypergraph $H=(\{x,y\},\{\{x,y\}\})$ gives:
$$G(H)=\langle x,y,e_{xy} \mid [[x,y],x], [[x,y],y], [x,y]=e_{xy} \rangle \cong \mathbb{H}$$
the Heisenberg group.
\end{example}

\end{section}

\begin{section}{Commutator Conventions}
Before proceeding, we establish conventions for manipulating commutators in 2-nilpotent contexts.

\begin{remark}\label{rem:commutator_identities}
The following identities hold in any 2-nilpotent group, and in particular
in $\langle e \rangle$ for any $e \in E$ and in the Heisenberg group of
Example~1.10:

    \begin{align*}
    [a,b] &= a^{-1}b^{-1}ab \\
    [a^{-1},b] &= ab^{-1}a^{-1}b = [a,b]^{-1}\\
    [a,b^{-1}] &= a^{-1}bab^{-1} = [a,b]^{-1}\\
    [b,a] &= [a,b]^{-1}\\
    [a^{-1},b^{-1}] &= [a,b]
    \end{align*}

    Thus any commutator of letters $a,b$ can be normalized to $[a,b]^{\pm 1}$.
\end{remark}

In our setup with ordered vertices $v_i < v_j$, we always express commutators as $e_{ij}^{\pm 1}$ where $i < j$.

Also in the Heisenberg group we have $[[a,b],a]=[[a,b],b]=1$, which with our extended alphabet means $[e_{ij},v_i]=[e_{ij},v_j]=1$.

\end{section}

\begin{section}{The Hyperalphabet and Monoid Structure}

\subsection{The Hyperalphabet $\mathcal{A}$}

\begin{definition}
A hyperedge $e \in E$ is \textbf{maximal} if there is no $e' \in E$ with
$e \subsetneq e'$. We denote by $\mathcal{E}_{\max}$ the set of all maximal
hyperedges of $\chi$. Note that under the density assumption
(Assumption~\ref{assump:dense}), every hyperedge is contained in at least
one maximal hyperedge, so every element of $A$ belongs to some
$\bar{E}$ with $E \in \mathcal{E}_{\max}$.
\end{definition}

\begin{definition}
For $u\subseteq A^{\pm}$, define $|u| = \{ a \mid a^{\epsilon} \in u \text{ for some } \epsilon \in \{\pm 1\}\}$.

The \textbf{hyperalphabet} is:
$\mathcal{A}=\{ [u] \mid u\subseteq A^{\pm}, |u|\subseteq \bar{E} \text{ for some } E \in \mathcal{E}_{\max}, \; x^{\epsilon}\in u \implies x^{-\epsilon}\not\in u \}$

Elements of $\mathcal{A}$ are called \textbf{hyper-letters}. Note that
$[\emptyset] \in \mathcal{A}$ is admitted (the conditions hold vacuously)
and is used in rule R3.
\end{definition}

\begin{remark}[Size of the Hyperalphabet]
For a hyperedge $E \in \mathcal{E}_{\max}$, the extended edge $\bar{E}$ has size $|E| + \binom{|E|}{2}$. This gives at most $2^{|\bar{E}|} - 1$ non-empty hyper-letters associated to $E$.
\end{remark}

\subsection{Canonical Maps Between Monoids}

\begin{definition}
Define the \textbf{interpretation map} $\mathcal{F}: \mathcal{A} \to (A^{\pm})^*$ by:
$$\mathcal{F}([u]) = a_1^{\epsilon_1} a_2^{\epsilon_2} \cdots a_k^{\epsilon_k}$$
where $u = \{a_1^{\epsilon_1}, \ldots, a_k^{\epsilon_k}\}$ and $a_1 < a_2 < \cdots < a_k$ in the total order on $A$. In particular, $\mathcal{F}([\emptyset])$ is the empty word. This extends to a monoid homomorphism $\mathcal{F}: \mathcal{A}^* \to (A^{\pm})^*$ by $\mathcal{F}([u_1]\cdots[u_k]) = \mathcal{F}([u_1])\cdots\mathcal{F}([u_k])$.
\end{definition}

\begin{example}
Let $u = \{v_3, v_5^{-1}, e_{3,5}^{-1}\}$ and $u_2 = \{v_2, v_4^{-1}\}$. Then:
\begin{itemize}
    \item $\mathcal{F}([u]) = v_3 v_5^{-1} e_{3,5}^{-1}$
    \item $\mathcal{F}([u_2]) = v_2 v_4^{-1}$
    \item $\mathcal{F}([u][u_2]) = v_3 v_5^{-1} e_{3,5}^{-1} v_2 v_4^{-1}$
\end{itemize}
\end{example}

\begin{proposition}
The map $\mathcal{F}: \mathcal{A}^* \to (A^{\pm})^*$ is surjective.
\end{proposition}
\begin{proof}
Since every $a \in A$ belongs to some $E \in \mathcal{E}_{\max}$, we have $[\{a^{\epsilon}\}] \in \mathcal{A}$ for all $a^{\epsilon} \in A^{\pm}$. Given $w = a_1^{\epsilon_1} \cdots a_n^{\epsilon_n}$, we have $w = \mathcal{F}([\{a_1^{\epsilon_1}\}] \cdots [\{a_n^{\epsilon_n}\}])$.
\end{proof}

\begin{remark}
    $\mathcal{F}$ is not injective. For example take the following two elements in $\mathcal{A}^*$: the empty hyperletter $[]$ and the empty word $\lambda$. They are not the same element of $\mathcal{A}^*$ since the empty hyperletter $[]$ is a hyperword of length 1, and $\lambda$ is of length 0, but both get mapped to the empty word on $(A^{\pm})^*$
\end{remark}

\begin{definition}
Define the \textbf{lifting map} $\sigma: (A^{\pm})^* \to \mathcal{A}^*$ by:
$$\sigma(a_1^{\epsilon_1} \cdots a_n^{\epsilon_n}) = [\{a_1^{\epsilon_1}\}] \cdots [\{a_n^{\epsilon_n}\}]$$
\end{definition}

Clearly $\mathcal{F} \circ \sigma = \text{Id}_{(A^{\pm})^*}$. Let $\mu: (A^{\pm})^* \to G$ denote the canonical quotient map.

Now we have three levels to work on:
\begin{enumerate}
    \item The group $G$ 
    \item The monoid of words $(A^{\pm})^*$
    \item The monoid of hyperwords $\mathcal{A}^*$
\end{enumerate}

In the group we have elements, in the first monoid words, and in the second one hyperwords. The idea now will be to lift problems in the group, like the word problems, to the word monoid, and then to the hyperword monoid, which has a lot more elements, but a structure that is easier to work with.

\end{section}

\begin{section}{The Rewriting System}

\subsection{The Price Functions}

\begin{definition}[Addability]
Let $[u] \in \mathcal{A}$ and $b^{\eta} \in A^{\pm}$. We say $b^{\eta}$ is \textbf{addable into $u$} if:
\begin{enumerate}
    \item $b \notin |u|$
    \item There exists $E \in \mathcal{E}_{\max}$ with $|u| \cup \{b\} \subseteq \bar{E}$
\end{enumerate}

When $b^{\eta}$ is addable into $u$, we write $[u + b^{\eta}]$ for the hyper-letter $u \cup \{b^{\eta}\}$.
\end{definition}

\begin{definition}[C function]
    Let $a,b\in A, a<b$ with $a$ and $b$ sharing an edge. Define $\mathbf{c}:A\times A \rightarrow A^{\pm}$ by:
    
    $$\mathbf{c}(a,b)=\begin{cases}
        e_{a,b} & a,b\in V
        \\
        - & \text{either }a\in E_2 \text{ or } b\in E_2
    \end{cases}$$

    and $\mathbf{c}(b,a)= \mathbf{c}(a,b)^{-1} \in A^{\pm}$.

    This function simulates the behaviour of the group relation of commutators at the level of the symbols. It can be thought of as the abbreviation (after group relations) of $[a,b]$, or equivalently, as the preimage of that group element.
\end{definition}


\begin{definition}[Price Functions]\label{def:price}
Let $[u] \in \mathcal{A}$ where $u = \{a_1^{\varepsilon_1}, \ldots, a_r^{\varepsilon_r}, c_1^{\delta_1}, \ldots, c_s^{\delta_s}\}$ with $a_1 < \cdots < a_r < b < c_1 < \cdots < c_s$ (note $b \notin |u|$ by addability). Suppose $b^{\eta}$ is addable into $u$.

Define the \textbf{left price}:
$$\mathbf{l_{b^{\eta}}}([u]) = [\{\mathbf{c}(a_i, b)^{\eta \varepsilon_i} \mid i=1,\ldots,r \}]$$

Define the \textbf{right price}:
$$\mathbf{r_{b^{\eta}}}([u]) = [\{\mathbf{c}(c_j, b)^{\eta \delta_j} \mid j=1,\ldots,s \}]$$

If $b^{\eta}$ is not addable into $u$, these functions are undefined.
\end{definition}

\begin{remark}
Price functions are edge-only hyperletters, which count the prices paid for interactions between a letter $b$ and a hyperletter $[u]$, assuming that $b$ is addable into $[u]$ then two interactions can happen:
\begin{enumerate}
    \item $b$ gets added into $[u]$ from the right. Then to order $[u]$ $b$ has to \textit{jump} through all letters that are larger than it. At each jump it "pays the price" of the corresponding commutator. Then, at the group level, all commutators can be grouped at the right since they commute with all other letters in $[u]$. All those commutators correspond to the \textbf{right price} $\mathbf{r_{b}}[u]$.
    \item $b$ gets removed from $[u]$ to the left. Then $b$ has to \textit{jump} through all letters that are smaller than it. At each jump it "pays the price" of the corresponding commutator. Then, at the group level, all commutators can be grouped at the left since they commute with all other letters in $[u]$. All those commutators correspond to the \textbf{left price} $\mathbf{l_{b}}[u]$.
\end{enumerate}
\end{remark}

\begin{lemma}\label{lem:price_welldef}
If $b^{\eta}$ is addable into $[u]$, then $\mathbf{l_{b^{\eta}}}([u])$ and $\mathbf{r_{b^{\eta}}}([u])$ are well-defined hyper-letters in $\mathcal{A}$. Moreover, both price hyper-letters contain only elements of $E_2^{\pm}$.
\end{lemma}

\begin{proof}
Since $b^\eta$ is addable into $u$, there exists $E \in \mathcal{E}_{\max}$ with $|u| \cup \{b\} \subseteq \bar{E}$.

For the left price $\mathbf{l_{b^{\eta}}}([u])$:
\begin{itemize}
    \item Each $a_i \in |u| \subseteq \bar{E}$ and $b \in \bar{E}$. If $a_i, b \in V$, then by density (Assumption~\ref{assump:dense}), $\{a_i, b\} \in E_2$, so $\mathbf{c}(a_i,b)= e_{a_i b} \in \overline{E} \subseteq A$.
    \item If $a_i \in E_2$ (or $b \in E_2$), then $\mathbf{c}(a_i,b) = -$ by Definition~4.3 (second case), so this term contributes nothing.
    \item The surviving commutators $[a_i, b]^{\eta\varepsilon_i}$ all lie in $A^\pm$, and since they are commutators of elements of $\bar{E}$, the set $\{[a_i,b]\}$ is contained in $E_2 \subseteq \bar{E}$. Thus $\mathbf{l_{b^\eta}}([u]) \in \mathcal{A}$ and contains only elements of $E_2^\pm$.
\end{itemize}

The argument for $\mathbf{r_{b^{\eta}}}([u])$ is identical, replacing $a_i$ with $c_j$.

Finally, we verify that both price hyper-letters satisfy the no-inverse-pair
condition of Definition~3.2. For the left price: the contributing elements
are $[a_i, b]^{\eta\varepsilon_i}$ for distinct $a_i < b$, giving distinct
commutator generators $e_{a_i b}$, so no generator appears with both
exponents. For the right price: the contributing elements are
$[c_j, b]^{\eta\delta_j}$ for distinct $c_j > b$, giving distinct generators
$e_{bc_j}$, again with no repeated generator. 
Hence both $\mathbf{l_{b^\eta}}([u])$ and $\mathbf{r_{b^\eta}}([u])$ are
well-formed hyper-letters in $\mathcal{A}$.
\end{proof}

\begin{remark}[Notational conventions for price properties]
We use two distinct operations which must not be conflated:
\begin{itemize}
    \item \textbf{Set union in a hyper-letter:} For hyper-letters $[u],[v]\in\mathcal{A}$
    with $|u|\cap|v|=\emptyset$ and $|u|\cup|v|\subseteq\bar{E}$ for some
    $E\in\mathcal{E}_{\max}$, the notation $[u+v]$ denotes the single
    hyper-letter $[u\cup v]\in\mathcal{A}$.
    \item \textbf{Concatenation in $\mathcal{A}^*$:} The product $[u]\cdot[v]$
    denotes the concatenation of two hyper-letters as a word of length 2
    in $\mathcal{A}^*$.
\end{itemize}

In the properties below, $\sim$ denotes rewriting equivalence under $\to^*$.
In Property~1, the left side $\mathbf{r_a}[u+v]$ is the $a$-price of the
single hyper-letter $[u\cup v]$, while the right side
$\mathbf{r_a}[u]\cdot\mathbf{r_a}[v]$ is their concatenation in $\mathcal{A}^*$.

Note also that Properties~5 and~6 handle both orderings of any two elements
uniformly: if $x \leq a$ then $x$ contributes nothing to the right $a$-price
(Property~5), and if $x \geq a$ then $x$ contributes nothing to the left
price (Property~6). In particular, no case split on the relative order of
two addable elements is needed when applying these properties.
\end{remark}

\begin{lemma}\label{lem:price_properties}
Whenever the involved symbols are well defined the price functions satisfy:
\begin{enumerate}
    \item $\mathbf{r_a}[u+v] \sim \mathbf{r_a}[u] \cdot \mathbf{r_a}[v]$ and $\mathbf{l_a}[u+v] \sim \mathbf{l_a}[u] \cdot \mathbf{l_a}[v]$ (distributivity)
    \item $\mathbf{r_{x^\eta}}[\{y^\delta\}] \cdot \mathbf{l_{y^\delta}}[\{x^\eta\}] \sim \lambda$ for all $\eta,\delta\in\{\pm 1\}$ (cancellation)
    \item $\mathbf{r_a}[u] \cdot \mathbf{r_b}[u] \sim \mathbf{r_b}[u] \cdot \mathbf{r_a}[u]$ (commutativity of prices)
    \item $[u] \cdot \mathbf{r_a}[u] \sim \mathbf{r_a}[u] \cdot [u]$ (prices commute with original)
    \item $\mathbf{r_a}[u+x] = \mathbf{r_a}[u]$ if $x\leq a$
    \item $\mathbf{l_a}[u+x] = \mathbf{l_a}[u]$ if $x\geq a$
    \item $\mathbf{r_a}[u] \cdot \mathbf{r_{a^{-1}}}[u] \sim \lambda$ and $\mathbf{l_a}[u] \cdot \mathbf{l_{a^{-1}}}[u] \sim \lambda$ (inverse cancellation)
\end{enumerate}
\end{lemma}

\begin{proof}
\textbf{Property 1: Distributivity.}

We prove $r_a[u + v] \sim r_a[u] \cdot r_a[v]$; the proof for $l$ is identical.

Since $[u + v] \in \mathcal{A}$ is a well-formed hyper-letter, we have
$|u| \cap |v| = \emptyset$ and $|u| \cup |v| \subseteq \bar{E}$ for some
$E \in \mathcal{E}_{\max}$, with $a \in \bar{E}$ by addability.
Partition the elements of $u$ and $v$ at $a$:
\begin{align*}
  u_R &= \{ b^{\varepsilon} \in u \mid b > a \}, \\
  v_R &= \{ b^{\varepsilon} \in v \mid b > a \}.
\end{align*}
By Definition~\ref{def:price}, the right $a$-prices are:
\begin{align*}
  r_a[u]   &= \bigl[\{ \mathbf{c}(b,a)^{\eta\varepsilon} \mid b^{\varepsilon} \in u_R \}\bigr], \\
  r_a[v]   &= \bigl[\{ \mathbf{c}(b,a)^{\eta\varepsilon} \mid b^{\varepsilon} \in v_R \}\bigr], \\
  r_a[u+v] &= \bigl[\{ \mathbf{c}(b,a)^{\eta\varepsilon} \mid b^{\varepsilon} \in u_R \cup v_R \}\bigr].
\end{align*}

\noindent\textit{(i) Disjoint support.}
Since $|u| \cap |v| = \emptyset$, we have $u_R \cap v_R = \emptyset$, so
$|r_a[u]| \cap |r_a[v]| = \emptyset$.

\noindent\textit{(ii) Common hyperedge.}
By Lemma~\ref{lem:price_welldef}, all elements of $r_a[u]$ and $r_a[v]$ lie in
$E_2^{\pm} \subseteq \bar{E}^{\pm}$, so $|r_a[u]| \cup |r_a[v]| \subseteq \bar{E}$.

\noindent\textit{(iii) Addability.}
By (i) and (ii), every generator $e^{\delta} \in r_a[v]$ satisfies
$e \notin |r_a[u]|$ and $|r_a[u]| \cup \{e\} \subseteq \bar{E}$,
so $e^{\delta}$ is addable into $r_a[u]$.

We now show $r_a[u] \cdot r_a[v] \to^* r_a[u+v]$ by
applying R1 to move each element of $r_a[v]$ into $r_a[u]$ one at a time.
The price terms arising at each step are trivial: every element of $r_a[u]$
and $r_a[v]$ lies in $E_2^{\pm}$, and $\mathbf{c}$ of two $E_2$-elements equals $-$ by Definition~4.3, so all prices are $[\emptyset]$. After applying R3 to delete each $[\emptyset]$ factor,
iterating over all $|v_R|$ elements of $r_a[v]$ (addability is preserved at
each step since disjointness and common hyperedge are maintained by (i)--(iii)):
\[
  r_a[u] \cdot r_a[v]
  \;\to^*\;
  \bigl[\{ \mathbf{c}(b,a)^{\eta\varepsilon} \mid b^{\varepsilon} \in u_R \cup v_R \}\bigr]
  = r_a[u + v].
\]
The proof for $l_a[u+v] \sim l_a[u] \cdot l_a[v]$ is identical.

\textbf{Property 2: Cancellation (general exponents).}

Let $\eta,\delta \in \{\pm 1\}$ and assume $x < y$.
By Definition~\ref{def:price}:
\begin{align*}
\mathbf{r_{x^\eta}}[\{y^\delta\}] &= [\{\mathbf{c}(y,x)^{\eta\delta}\}] = [\{e_{xy}^{-\eta\delta}\}]\\
\mathbf{l_{y^\delta}}[\{x^\eta\}] &= [\{\mathbf{c}(x,y)^{\eta\delta}\}] = [\{e_{xy}^{\eta\delta}\}]
\end{align*}
Concatenating: $[\{e_{xy}^{-\eta\delta}\}][\{e_{xy}^{\eta\delta}\}]
\xrightarrow{R_2} [\emptyset][\emptyset] \xrightarrow{R_3} [\emptyset]\xrightarrow{R_3}  \lambda $.
The case $x > y$ is analogous with $e_{yx}$ in place of $e_{xy}$.

\textbf{Property 3: Commutativity of prices.}

Both $\mathbf{r_a}[u]$ and $\mathbf{r_b}[u]$ consist of edge hyperletters. Since $\mathbf{c}$ of two $E_2$-elements equals $-$ by Definition~4.3, moving any element of $\mathbf{r_b}[u]$ past any element of $\mathbf{r_a}[u]$ via R1 produces only $[\emptyset]$ price terms, which vanish by R3. Hence the two hyperletters can be freely reordered.

\textbf{Property 4: Prices commute with original.}

$\mathbf{r_a}[u]$ consists of edge hyperletters $e_{c_j,a} \in E_2^\pm$. Since $\mathbf{c}$ of an $E_2$-element with any other element equals $-$ by Definition~4.3, moving any element of $\mathbf{r_a}[u]$ past any element of $[u]$ via R1 produces only $[\emptyset]$ price terms. Hence $[u]$ and $\mathbf{r_a}[u]$ commute under $\to^*$.

\textbf{Property 5: $\mathbf{r_a}[u+x] = \mathbf{r_a}[u]$ if $x \leq a$.}

The right price involves only elements of $u$ greater than $a$. Since $x \leq a$, $x$ does not contribute:
$$\mathbf{r_a}[u+x] = [\{\mathbf{c}(w,a)^{\eta\epsilon_w} \mid w \in (u \cup \{x\}), w > a\}] = [\{\mathbf{c}(w,a)^{\eta\epsilon_w} \mid w \in u, w > a\}] = \mathbf{r_a}[u]$$

\textbf{Property 6: $\mathbf{l_a}[u+x] = \mathbf{l_a}[u]$ if $x \geq a$.}

Analogous: the left price involves only elements less than $a$, and $x \geq a$ does not contribute.

\textbf{Property 7: Inverse cancellation.}

The elements of $u$ greater than $a$ are $c_{i_1}, \ldots, c_{i_m}$.

\begin{align*}
\mathbf{r_a}[u] &= [\{\mathbf{c}(c_{i_j}, a)^{\epsilon_{i_j}}\}_{j=1}^m]\\
\mathbf{r_{a^{-1}}}[u] &= [\{\mathbf{c}(c_{i_j}, a^{-1})^{\epsilon_{i_j}}\}_{j=1}^m] = [\{\mathbf{c}(c_{i_j}, a)^{-\epsilon_{i_j}}\}_{j=1}^m]
\end{align*}
where we used $\mathbf{c}(c, a^{-1}) = \mathbf{c}(c,a)^{-1}$ (Remark~\ref{rem:commutator_identities}). Applying R2 and R3 repeatedly:
$$\mathbf{r_a}[u] \cdot \mathbf{r_{a^{-1}}}[u] \xrightarrow{R_2^m, R_3} [\emptyset] \xrightarrow{R_3}\lambda$$
The proof for $\mathbf{l}$ is identical.
\end{proof}

\begin{lemma}[Interpretation in the Group]\label{lem:interpretation}
In the group $G$, the price functions satisfy the following identities for
all $[u] \in \mathcal{A}$ and all $b^\eta$ addable into $u$:
\begin{align}
\mathcal{F}([u + b^{\eta}]) &= b^{\eta} \cdot \mathcal{F}([u]) \cdot \mathcal{F}(\mathbf{l_{b^{\eta}}}([u])) \label{eq:left_id}\\
\mathcal{F}([u]) \cdot b^{\eta} &= \mathcal{F}([u + b^{\eta}]) \cdot \mathcal{F}(\mathbf{r_{b^{\eta}}}([u])) \label{eq:right_id}
\end{align}
That is, $\mathbf{l}$ gives the commutators generated when moving $b^{\eta}$
to the left past $u$, and $\mathbf{r}$ gives those generated when moving
$b^\eta$ into $[u+b^\eta]$ from the right.
\end{lemma}

\begin{proof}
We prove \eqref{eq:right_id} by induction on $|u|$; the proof of
\eqref{eq:left_id} is analogous.

\textit{Base case} ($|u|=1$): Write $u = \{c^{\delta}\}$. If $c < b$
then $c$ is in the left part and the right price is empty:
$\mathbf{r_{b^\eta}}([u]) = [\emptyset]$, and $\mathcal{F}([u]) \cdot b^\eta
= c^\delta b^\eta = b^\eta c^\delta [\emptyset] = \mathcal{F}([u+b^\eta])
\cdot \mathcal{F}([\emptyset])$ after noting $\mathcal{F}([\emptyset])=1$
in the group.

If $c > b$, using the identity $xy = yx[x,y]$ in the group:
$$c^\delta b^\eta = b^\eta c^\delta [c^\delta, b^\eta] = b^\eta c^\delta [c,b]^{\delta\eta}$$
Since $[c,b]^{\delta\eta} \in E_2^\pm$ by Lemma~\ref{lem:price_welldef}, and
$\mathbf{r_{b^\eta}}([\{c^\delta\}]) = [\{[c,b]^{\delta\eta}\}]$, we have
$\mathcal{F}([u])\cdot b^\eta = \mathcal{F}([u+b^\eta])\cdot
\mathcal{F}(\mathbf{r_{b^\eta}}([u]))$ as required.

\textit{Inductive step}: Suppose \eqref{eq:right_id} holds for all hyper-letters
of size $< |u|$. We split $u = u_L \sqcup u_R$
where $u_L = \{a_i^{\varepsilon_i} \in u \mid a_i < b\}$ and
$u_R = \{c_j^{\delta_j} \in u \mid c_j > b\}$, so $u = u_L \cup u_R$.

In the group, $\mathcal{F}([u]) = \mathcal{F}([u_L])\cdot\mathcal{F}([u_R])$
since $\mathcal{F}$ outputs elements in sorted order and all elements of $u_L$
precede all elements of $u_R$ in the total order.

By induction applied to $u_R$ (which has size $< |u|$):
$$\mathcal{F}([u_R])\cdot b^\eta = \mathcal{F}([u_R + b^\eta])\cdot
\mathcal{F}(\mathbf{r_{b^\eta}}([u_R]))$$

Since elements of $u_L$ are all less than $b^\eta$ and commutators of $E_2$-type
are central by 2-nilpotency:
$$\mathcal{F}([u_L])\cdot\mathcal{F}([u_R + b^\eta]) = \mathcal{F}([u_L + u_R + b^\eta])
= \mathcal{F}([u + b^\eta])$$

Also, $\mathcal{F}([u_L])$ commutes with $\mathcal{F}(\mathbf{r_{b^\eta}}([u_R]))$:
the latter consists of commutators in $E_2^\pm$ which are central in $\langle e\rangle$
by 2-nilpotency, and elements of $u_L$ lie in the same hyperedge. Thus:
$$\mathcal{F}([u])\cdot b^\eta = \mathcal{F}([u_L])\cdot\mathcal{F}([u_R])\cdot b^\eta
= \mathcal{F}([u+b^\eta])\cdot\mathcal{F}(\mathbf{r_{b^\eta}}([u_R]))$$

$\mathbf{r_{b^\eta}}([u_R]) =
\mathbf{r_{b^\eta}}([u])$ by Property~5 of Lemma~\ref{lem:price_properties}, since $u_L$ consists entirely of elements $\leq b$ and thus contributes nothing to the right price.
This completes the induction.
\end{proof}

\begin{example}
Let $\chi$ be a hypergraph containing the hyperedge $\{v_1, v_2, v_3\} \in E_3$,
so that by density $\{v_1,v_2\}, \{v_1,v_3\}, \{v_2,v_3\} \in E_2$.
Let $u = \{v_1, v_3^{-1}, e_{13}\}$ and add $v_2$ (with $v_1 < v_2 < v_3$, $\eta=1$). Elements of $u$ less than $v_2$: $\{v_1\}$. Elements greater than $v_2$: $\{v_3^{-1}\}$ (and $e_{13}$, but $[v_2, e_{13}]=1$ by 2-nilpotency since $\{v_1,v_2,v_3\}\in E_3$). Then:
\begin{align*}
\mathbf{l_{v_2}}([u]) &= [\{[v_1, v_2]\}] = [\{e_{12}\}]\\
\mathbf{r_{v_2}}([u]) &= [\{[v_3, v_2]^{-1}\}] = [\{e_{23}\}]
\end{align*}
(since $[v_3, v_2] = e_{23}^{-1}$, so $[v_3, v_2]^{-1} = e_{23}$).
\end{example}

\subsection{The Rewriting Rules}

\begin{definition}[Rewriting System]\label{def:rewriting}
The rewriting system on $\mathcal{A}^*$ consists of three rules:

\textbf{(R1) Movement:}
$$w_1[u][v+x^{\epsilon}]w_2\rightarrow w_1[u+x^{\epsilon}]\mathbf{r_{x^{\epsilon}}}([u])[v]\mathbf{l_{x^{\epsilon}}}([v])w_2$$
provided $x^{\epsilon}$ is addable into $u$ and $v$.

\textbf{(R2) Cancellation:}
$$w_1[u+x^{-\epsilon}][v+x^{\epsilon}]w_2\rightarrow w_1[u]\mathbf{r_{x^{\epsilon}}}([u])[v]\mathbf{l_{x^{\epsilon}}}([v])w_2$$
provided $x^{\epsilon}$ is addable into $u$ and $v$.

\textbf{(R3) Empty Deletion:}
$$w_1[\emptyset]w_2\rightarrow w_1w_2$$
\end{definition}

\begin{definition}
A word $w \in \mathcal{A}^*$ is \textbf{reduced} (in \textbf{normal form}) if no rewriting rule applies.
\end{definition}

\end{section}

\begin{section}{Confluence and Termination}

\begin{newmanslemma}[\cite{Newman42}]
    A terminating rewriting system is confluent if and only if it is locally confluent.
\end{newmanslemma}

\subsection{Local Confluence}

\begin{lemma}\label{lem:local_confluence}
    The rewriting system is locally confluent.
\end{lemma}

\begin{proof}
We consider all cases where two rules apply to the same word. We write $R_i$ for single rule applications and $P_i$ for use of property $i$ of Lemma~\ref{lem:price_properties}.

\textbf{Case 0: Non-overlapping reductions.} Immediate --- the two reductions act on disjoint pairs of hyper-letters and commute.

\textbf{Case 1: One reduction involves R3.}

Subcase (1.1): $R$ is R1, $R'$ is R3, applied to $w = w_1[u][\emptyset][v+x^\epsilon]w_2$:

\vspace{20pt}
\begin{center}
\begin{adjustbox}{center,margin=0pt -1.5cm}
\begin{tikzcd}[column sep=small]
w_1[u][\emptyset][v+x^\epsilon]w_2 \arrow[r,"R_3"] \arrow[d,"R_1"] & w_1[u][v+x^\epsilon]w_2 \arrow[r,"R_1"]&  w_1[u+x^\epsilon]\mathbf{r_{x^\epsilon}}[u][v]\mathbf{l_{x^\epsilon}}[v]w_2 \\
w_1[u][\emptyset+x^\epsilon]\mathbf{r_{x^\epsilon}}[\emptyset][v]\mathbf{l_{x^\epsilon}}[v]w_2 \arrow[r,"R_1"] & w_1[u+x^\epsilon]\mathbf{r_{x^\epsilon}}[u][\emptyset]\mathbf{l_{x^\epsilon}}[\emptyset]\mathbf{r_{x^\epsilon}}[\emptyset][v]\mathbf{l_{x^\epsilon}}[v]w_2 \arrow[r,"R_3^3"] & w_1[u+x^\epsilon]\mathbf{r_{x^\epsilon}}[u][v]\mathbf{l_{x^\epsilon}}[v]w_2  \arrow[u] 
\end{tikzcd}
\end{adjustbox}
\end{center}
\vspace{30pt}

Subcase (1.2) for R2 and R3 is similar.

\textbf{Case 2: Both R1.}

In this case we assume:
\begin{itemize}
  \item $y^\delta$ is addable into $u$ (witnessing $E_1 \in \mathcal{E}_{\max}$
        with $|u| \cup \{y\} \subseteq \bar{E}_1$), and
  \item $x^\epsilon$ is addable into $v+y^\delta$ (witnessing $E_2 \in
        \mathcal{E}_{\max}$ with $|v| \cup \{x,y\} \subseteq \bar{E}_2$).
\end{itemize}
These are precisely the conditions under which both R1 rules apply
simultaneously to $w = w_1[u][v+y^\delta][w+x^\epsilon]w_2$. Note that since
$|v|\cup\{x\} \subseteq |v|\cup\{x,y\} \subseteq \bar{E}_2$, the element $x^\epsilon$
is also addable into $v$, as required for the lower path.

\vspace{40pt}
\begin{center}
\begin{adjustbox}{center,margin=0pt -1.5cm}
\begin{tikzcd}[column sep=tiny, row sep=large]
w_1[u+y^{\delta}]\mathbf{r_{y^{\delta}}}[u][v+x^{\epsilon}]\mathbf{l_{y^{\delta}}}[v+x^\epsilon]\mathbf{r_{x^{\epsilon}}}[v+y^\delta][w]\mathbf{l_{x^{\epsilon}}}[w]w_2 \arrow[r,"P_1"] & w_1[u+y^{\delta}]\mathbf{r_{y^{\delta}}}[u][v+x^{\epsilon}]\mathbf{l_{y^{\delta}}}[v]\mathbf{l_{y^{\delta}}}[x^\epsilon]\mathbf{r_{x^{\epsilon}}}[y^\delta]\mathbf{r_{x^{\epsilon}}}[v][w]\mathbf{l_{x^{\epsilon}}}[w]w_2 \arrow[d,"P_3\,P_2"] \\
w_1[u][v+x^\epsilon+y^\delta]\mathbf{r_{x^\epsilon}}[v+y^\delta][w]\mathbf{l_{x^\epsilon}}[w]w_2 \arrow[u,"R_1(y)"] &  w_1[u+y^{\delta}]\mathbf{r_{y^{\delta}}}[u][v+x^{\epsilon}]\mathbf{l_{y^{\delta}}}[v][\emptyset]\mathbf{r_{x^{\epsilon}}}[v][w]\mathbf{l_{x^{\epsilon}}}[w]w_2 \arrow[d,"R_3"]\\
w_1[u][v+y^\delta][w+x^\epsilon]w_2 \arrow[u,"R_1(x)"] \arrow[d,"R_1(y)"] & w_1[u+y^{\delta}]\mathbf{r_{y^{\delta}}}[u][v+x^{\epsilon}]\mathbf{l_{y^{\delta}}}[v]\mathbf{r_{x^{\epsilon}}}[v][w]\mathbf{l_{x^{\epsilon}}}[w]w_2  \\
w_1[u+y^\delta]\mathbf{r_{y^\delta}}[u][v]\mathbf{l_{y^{\delta}}}[v][w+x^\epsilon]w_2 \arrow[r,"R_1(x)"] & w_1[u+y^{\delta}]\mathbf{r_{y^\delta}}[u][v+x^\epsilon]\mathbf{r_{x^\epsilon}}[v]\mathbf{l_{y^{\delta}}}[v][w]\mathbf{l_{x^\epsilon}}[w]w_2 \arrow[u,"P_3"]
\end{tikzcd}
\end{adjustbox}
\end{center}
\vspace{40pt}

The step labeled $P_3\,P_2$ first applies $P_3$ (commutativity of prices) to swap $\mathbf{l_{y^\delta}}[x^\epsilon]$ and $\mathbf{r_{x^\epsilon}}[y^\delta]$ into the order $\mathbf{r_{x^\epsilon}}[y^\delta]\cdot\mathbf{l_{y^\delta}}[x^\epsilon]$, and then applies $P_2$ (cancellation) to reduce $\mathbf{r_{x^\epsilon}}[\{y^\delta\}]\cdot\mathbf{l_{y^\delta}}[\{x^\epsilon\}]\sim\lambda$.

\textbf{Case 3: Both R2.}


Subcase (3.1): Different symbols $x \neq y$ in $w = w_1[u+x^{-\epsilon}][v+x^\epsilon+y^{-\delta}][w+y^\delta]w_2$:
\vspace{40pt}
\begin{center}
\begin{adjustbox}{center,margin=0pt -1.5cm}
\begin{tikzcd}[column sep=tiny, row sep=large]
w_1[u]\mathbf{r_{x^{\epsilon}}}[u][v]\mathbf{r_{y^{\delta}}}[v]\mathbf{l_{x^{\epsilon}}}[v+y^{-\delta}][w]\mathbf{l_{y^{\delta}}}[w]w_2 \arrow[r,"P_1"] & w_1[u]\mathbf{r_{x^{\epsilon}}}[u][v]\mathbf{r_{y^{\delta}}}[v]\mathbf{l_{x^{\epsilon}}}[v]\mathbf{l_{x^{\epsilon}}}[y^{-\delta}][w]\mathbf{l_{y^{\delta}}}[w]w_2 \arrow[d,"P_2"] \\
w_1[u]\mathbf{r_{x^{\epsilon}}} [u] [v+y^{-\delta}]\mathbf{l_{x^{\epsilon}}}[v+y^{-\delta}][w+y^\delta]w_2 \arrow[u,"R_2(y)"] & w_1[u]\mathbf{r_{x^{\epsilon}}}[u][v]\mathbf{l_{x^{\epsilon}}}[v]\mathbf{r_{y^{\delta}}}[v]\mathbf{l_{x^{\epsilon}}}[y^{-\delta}][w]\mathbf{l_{y^{\delta}}}[w]w_2  \arrow[d,"P_3"] \\
w_1[u+x^{-\epsilon}][v+x^\epsilon+y^{-\delta}][w+y^\delta]w_2 \arrow[u,"R_2(x)"] \arrow[d,"R_2(y)"] &  w_1[u]\mathbf{r_{x^{\epsilon}}}[u][v]\mathbf{l_{x^{\epsilon}}}[v]\mathbf{r_{y^{\delta}}}[v]\mathbf{r_{y^{\delta}}}[x^\epsilon][w]\mathbf{l_{y^{\delta}}}[w]w_2 \\
w_1[u+x^{-\epsilon}][v+x^\epsilon]\mathbf{r_{y^{\delta}}}[v+x^\epsilon][w]\mathbf{l_{y^{\delta}}}[w]w_2 \arrow[r,"R_2(x)"] & w_1[u]\mathbf{r_{x^{\epsilon}}}[u][v]\mathbf{l_{x^{\epsilon}}}[v]\mathbf{r_{y^{\delta}}}[v+x^\epsilon][w]\mathbf{l_{y^{\delta}}}[w]w_2 \arrow[u,"P_1"]
\end{tikzcd}
\end{adjustbox}
\end{center}
\vspace{40pt}

Subcase (3.2): Same element, $w = w_1[u+x^{-\epsilon}][v+x^\epsilon][w'+x^{-\epsilon}]w_2$.

Two rules overlap: R2 applies to pair $(1,2)$ cancelling $x^{-\epsilon}$ and $x^\epsilon$; R2 also applies to pair $(2,3)$ with $x^\epsilon$ in the left letter and $x^{-\epsilon}$ in the right letter.
\vspace{40pt}
\begin{center}
\begin{adjustbox}{center,margin=0pt -1.5cm}
\begin{tikzcd}[column sep=small]
w_1[u]\mathbf{r_{x^{\epsilon}}}[u][v]\mathbf{l_{x^\epsilon}}[v][w'+x^{-\epsilon}]w_2 \arrow[r,"R_2"] & w_1[u]\mathbf{r_{x^{\epsilon}}}[u][v+x^{-\epsilon}]\mathbf{l_{x^\epsilon}}[v]\mathbf{r_{x^{-\epsilon}}}[v][w']\mathbf{l_{x^{-\epsilon}}}[w']w_2 \arrow[d,"R_2"]   \\
w_1[u+x^{-\epsilon}][v+x^\epsilon][w'+x^{-\epsilon}]w_2 \arrow[u,"R_2\text{ on }(1-2)"] \arrow[d,"R_2\text{ on }(2-3)"]  & w_1[u+x^{-\epsilon}]\mathbf{r_{x^{-\epsilon}}}[u]\mathbf{r_{x^{\epsilon}}}[u][v]\mathbf{l_{x^{-\epsilon}}}[v]\mathbf{r_{x^{-\epsilon}}}[v]\mathbf{l_{x^\epsilon}}[v][w']\mathbf{l_{x^{-\epsilon}}}[w']w_2 \arrow[d,"P_5\,P_7"] \\
w_1[u+x^{-\epsilon}][v]\mathbf{r_{x^{-\epsilon}}}[v][w']\mathbf{l_{x^{-\epsilon}}}[w']w_2  & 
\arrow[l,"R_3"] w_1[u+x^{-\epsilon}][\emptyset][v]\mathbf{r_{x^{-\epsilon}}}[v][w']\mathbf{l_{x^{-\epsilon}}}[w']w_2
\end{tikzcd}
\end{adjustbox}
\end{center}
\vspace{40pt}

The step labeled $P_5\,P_7$ first applies $P_5$ to reduce $\mathbf{r_{x^\epsilon}}[u+x^{-\epsilon}] = \mathbf{r_{x^\epsilon}}[u]$ (since $x^{-\epsilon}$ and $x^\epsilon$ share the same base element, so $x^{-\epsilon} \leq x^\epsilon$ does not contribute to the right $x^\epsilon$-price), giving the factor $\mathbf{r_{x^{-\epsilon}}}[u]\cdot\mathbf{r_{x^\epsilon}}[u]$. Then $P_7$ (inverse cancellation) gives $\mathbf{r_{x^{-\epsilon}}}[u]\cdot\mathbf{r_{x^\epsilon}}[u] \sim \lambda$. The same reduction applies to the $\mathbf{l}$ factors.

\textbf{Case 4: R1 and R2.}

Subcase (4.1): Different elements, $w = w_1[u+x^{-\epsilon}][v+x^\epsilon][w'+y^\delta]w_2$:
\vspace{40pt}
\begin{center}
\begin{adjustbox}{center,margin=0pt -1.5cm}
\begin{tikzcd}[column sep=small, row sep=large]
w_1[u]\mathbf{r_{x^{\epsilon}}}[u][v]\mathbf{l_{x^{\epsilon}}}[v][w'+y^\delta]w_2 \arrow[r,"R_1(y)"] & w_1[u]\mathbf{r_{x^{\epsilon}}}[u][v+y^\delta]\mathbf{r_{y^\delta}}[v]\mathbf{l_{x^{\epsilon}}}[v][w']\mathbf{l_{y^\delta}}[w']w_2 \\
w_1[u+x^{-\epsilon}][v+x^\epsilon][w'+y^\delta]w_2 \arrow[u,"R_2(x)"] \arrow[d,"R_1(y)"] & w_1[u]\mathbf{r_{x^\epsilon}}[u][v+y^\delta]\mathbf{r_{y^\delta}}[v]\mathbf{r_{y^\delta}}[x^\epsilon]\mathbf{l_{x^\epsilon}}[y^\delta]\mathbf{l_{x^\epsilon}}[v][w']\mathbf{l_{y^\delta}}[w']w_2 \arrow[u,"P_6 P_2"] \\
w_1[u+x^{-\epsilon}][v+x^\epsilon+y^\delta]\mathbf{r_{y^\delta}}[v+x^\epsilon][w']\mathbf{l_{y^\delta}}[w']w_2 \arrow[r,"R_2(x)"] & w_1[u]\mathbf{r_{x^\epsilon}}[u][v+y^\delta]\mathbf{r_{y^\delta}}[v+x^\epsilon]\mathbf{l_{x^\epsilon}}[v+y^\delta][w']\mathbf{l_{y^\delta}}[w']w_2 \arrow[u,"P_5"]
\end{tikzcd}
\end{adjustbox}
\end{center}
\vspace{40pt}

Subcases (4.2) and (4.3) follow similar patterns.

\textbf{Case 5: Overlapping in a single hyper-letter.}

Subcase (5.1): $w = w_1[u][v+x^{\epsilon}+y^{\delta}]w_2$ where both $x,y$ are addable into $u$, $x < y$:

\vspace{40pt}
\begin{center}
\begin{adjustbox}{center,margin=0pt -1.5cm}
\begin{tikzcd}[column sep=tiny, row sep=large]
w_1[u+x^{\epsilon}+y^{\delta}]\mathbf{r_{y^{\delta}}}[u+x^{\epsilon}]\mathbf{r_{x^{\epsilon}}}[u][v]\mathbf{l_{y^{\delta}}}[v]\mathbf{l_{y^{\delta}}}[v+y^{\delta}]w_2 \arrow[r,"P_5 P_6"] & w_1[u+x^{\epsilon}+y^{\delta}]\mathbf{r_{y^{\delta}}}[u]\mathbf{r_{y^{\delta}}}[x^{\epsilon}]\mathbf{r_{x^{\epsilon}}}[u][v]\mathbf{l_{y^{\delta}}}[v]\mathbf{l_{y^{\delta}}}[v]w_2 \arrow[d,"P_3"] \\
w_1[u+x^{\epsilon}]\mathbf{r_{x^{\epsilon}}}[u][v+y^{\delta}]\mathbf{l_{y^{\delta}}}[v+y^{\delta}]w_2 \arrow[u,"R_1(y)"] & w_1[u+x^{\epsilon}+y^{\delta}]\mathbf{r_{x^{\epsilon}}}[u]\mathbf{r_{y^{\delta}}}[u]\mathbf{r_{y^{\delta}}}[x^{\epsilon}][v]\mathbf{l_{y^{\delta}}}[v]\mathbf{l_{y^{\delta}}}[v]w_2  \arrow[d,"P_3"]\\
w_1[u][v+x^{\epsilon}+y^{\delta}]w_2 \arrow[u,"R_1(x)"] \arrow[d,"R_1(y)"] & w_1[u+x^{\epsilon}+y^{\delta}]\mathbf{r_{x^{\epsilon}}}[u]\mathbf{r_{x^{\epsilon}}}[y^{\delta}]\mathbf{r_{y^{\delta}}}[u][v]\mathbf{l_{y^{\delta}}}[v]\mathbf{l_{y^{\delta}}}[v]w_2 \\
w_1[u+y^{\delta}]\mathbf{r_{y^{\delta}}}[u][v+x^{\epsilon}]\mathbf{l_{y^{\delta}}}[v+x^{\epsilon}]w_2 \arrow[r,"R_1(x)"] & w_1[u+x^{\epsilon}+y^{\delta}]\mathbf{r_{x^{\epsilon}}}[u+y^{\delta}]\mathbf{r_{y^{\delta}}}[u][v]\mathbf{l_{y^{\delta}}}[v]\mathbf{l_{y^{\delta}}}[v+x^{\epsilon}]w_2 \arrow[u,"P_5 P_6"]
\end{tikzcd}
\end{adjustbox}
\end{center}
\vspace{40pt}

Subcase (5.2) and (5.3) follow analogously using Lemma~\ref{lem:price_properties}.

In all cases we find a common descendant, proving local confluence.
\end{proof}

\subsection{Termination}

\begin{lemma}\label{lem:termination}
The rewriting system is terminating.
\end{lemma}

\begin{proof}
For a word $w = [u_1]\cdots[u_n] \in \mathcal{A}^*$, 

let $u_i^V = u_i \cap (V^{\pm})$ and $u_i^E = u_i \cap (E_2^{\pm})$, define the
\emph{vertex projection} $w|_V$ and \emph{edge projection} $w|_E$ as
the subsequences of $w$ consisting, respectively, of the
\emph{mixed} hyper-letters (those with $u_k^V \neq \emptyset$) and the
\emph{non-empty edge-only} hyper-letters (those with $u_k^V = \emptyset$,
$u_k \neq \emptyset$). Write:
\[
  w|_V = [v_1]\cdots[v_p], \qquad w|_E = [e_1]\cdots[e_q].
\]
Define the termination measure
\[
  f(w) \;=\;
  \Biggl(
    \underbrace{\sum_{k=1}^{p} k \cdot |v_k^V|}_{f_1},\;\;
    \underbrace{\sum_{k=1}^{q} k \cdot |e_k|}_{f_2},\;\;
    \underbrace{|w|}_{f_3}
  \Biggr)
  \;\in\; \mathbb{N}^3,
\]
ordered lexicographically. We show $f$ strictly decreases under each rule.

\medskip
\noindent\textbf{Rule R1.}
The rule moves $x^\varepsilon \in u_j$ into an earlier block $u_{j-1}$,
producing two price hyper-letters $r_{x^\varepsilon}([u_{j-1}])$ and
$l_{x^\varepsilon}([u_j])$ which, by Lemma~\ref{lem:price_welldef}, consist entirely of
elements of $E_2^\pm$.

\smallskip
\noindent\emph{Case (a): $x \in V$.}
The two price hyper-letters are pure-edge, so they do not appear in
$w|_V$. Hence the vertex projection sees only that $x$ moves from
mixed-rank $j$ to mixed-rank $j-1 < j$, with no new mixed blocks inserted.
The mixed-rank of every other block is unchanged. Therefore:
\[
  \Delta f_1 \;=\;  (j-1)-j = -1 \;<\; 0.
\]
The first component strictly decreases, so $f$ decreases regardless of
$f_2$ and $f_3$.

(Note: the insertion of the pure-edge price hyper-letters
$r_{x^\varepsilon}([u_{j-1}])$ and $l_{x^\varepsilon}([u_j])$ does not affect
$w|_V$, since those blocks do not appear in the vertex projection. Hence
the indices $j-1 < j$ refer to positions in $w|_V$ exclusively, and no
intermediate mixed-block index is disturbed. Furthermore, $f_3 = |w|$
strictly increases by $2$ under R1 since two hyper-letters are replaced by
four; however, since the lexicographic order on $\mathbb{N}^3$ is a
well-order and $\Delta f_1 = -1 < 0$, the strict decrease in $f_1$
dominates any increase in $f_3$.)

\smallskip
\noindent\emph{Case (b): $x \in E_2$.}
Since $x \in E_2^\pm$, Definition~4.3 gives $\mathbf{c}(x,y) = -$ for all $y$,
so both price hyper-letters equal $[\emptyset]$ by Definition~4.4 and do not appear in $w|_E$.
The edge projection sees only that $x$ moves from edge-rank $j$ to
edge-rank $j-1 < j$ (with all other edge-ranks unchanged), so:
\[
  \Delta f_2 \;=\; (j-1)-j \;=\; -1 \;<\; 0.
\]
The second component strictly decreases (with $f_1$ unchanged), so $f$
decreases.

\medskip
\noindent\textbf{Rule R2.}
Identical analysis to R1, with the additional removal of $x$ from both
blocks, giving a strictly greater decrease in $f_1$ or $f_2$ as
appropriate.

\medskip
\noindent\textbf{Rule R3.}
Removing $[\emptyset]$ does not affect $w|_V$ or $w|_E$ (the empty
block belongs to neither projection), so $f_1$ and $f_2$ are unchanged.
The word length $f_3 = |w|$ strictly decreases by $1$.

\medskip
\noindent Since $f$ takes values in $\mathbb{N}^3$ with the
lexicographic order (a well-order) and strictly decreases under every
rule, the rewriting system terminates.
\end{proof}

\begin{corollary}\label{cor:confluence}
The rewriting system is confluent.
\end{corollary}
\begin{proof}
By Newman's Lemma, since the system is terminating (Lemma~\ref{lem:termination}) and locally confluent (Lemma~\ref{lem:local_confluence}), it is confluent.
\end{proof}

\end{section}

\begin{section}{The Word Problem}

\subsection{Correctness of the Rewriting System}

\begin{lemma}
Every defining relator of
$G'(\chi)$ reduces to $[\emptyset]$ under $\to^*$.
\end{lemma}
\begin{proof}

We use the shorthand $[a^\epsilon]$ for the singleton hyper-letter
$[\{a^\epsilon\}]$, and write reductions in $\mathcal{A}^*$ after lifting
via $\sigma$.

The key observation used throughout is: for any $z$ and $e_{xy}$ lying
in the same extended hyperedge $\bar{E}$, we have $[z, e_{xy}] = 1$ by
the 2-nilpotent relations of $G'(\chi)$. Consequently, whenever $z$ or
$e_{xy}$ appears in a price computation involving such a pair, the
corresponding price is $[\emptyset]$.

\medskip
\noindent\textbf{Relator 1: $[[x,y],x] \to^* [\emptyset]$.}

In $G'(\chi)$, $[x,y]=e_{xy}$, so $[[x,y],x] = [e_{xy},x] = e_{xy}^{-1}x^{-1}e_{xy}x$.
Lifting: $[e_{xy}^{-1}][x^{-1}][e_{xy}][x]$, with $x < y < e_{xy}$.

\textit{Step 1.} Apply R1 to move $e_{xy}$ left past $[x^{-1}]$ (addable since $e_{xy} \notin |[x^{-1}]|$ and both lie in $\bar{E}$). The left price $\mathbf{l_{e_{xy}}}([x^{-1}])$ collects $[x, e_{xy}]^{-1} = 1$ (by the key observation). The right price is $[\emptyset]$ (no elements of $[x^{-1}]$ exceed $e_{xy}$). Thus:
$$[x^{-1}][e_{xy}] \to^* [x^{-1}+e_{xy}]$$

\textit{Step 2.} Apply R2 to $[e_{xy}^{-1}][x^{-1}+e_{xy}]$ (with $u=\emptyset$, $x^{-\epsilon}=e_{xy}^{-1}$, $v=\{x^{-1}\}$, $x^\epsilon=e_{xy}$). Prices: $\mathbf{r_{e_{xy}}}([\emptyset])=[\emptyset]$, $\mathbf{l_{e_{xy}}}([x^{-1}])=[\emptyset]$ as above. After R3:
$$[e_{xy}^{-1}][x^{-1}+e_{xy}] \to^* [x^{-1}]$$

\textit{Step 3.} Apply R2 to $[x^{-1}][x]$ with trivial prices:
$$[x^{-1}][x] \to^* [\emptyset]$$

\medskip
\noindent\textbf{Relator 2: $[[x,y],y] \to^* [\emptyset]$.}

Similarly, $[[x,y],y]=[e_{xy},y]=e_{xy}^{-1}y^{-1}e_{xy}y$.
Lifting: $[e_{xy}^{-1}][y^{-1}][e_{xy}][y]$, with $y < e_{xy}$.

Steps are identical to Relator~1 with $y$ in place of $x$:
\begin{itemize}
    \item R1 on $[y^{-1}][e_{xy}]$: left price $[y,e_{xy}]^{-1}=1$ (key observation), right price $[\emptyset]$. Gives $[y^{-1}+e_{xy}]$.
    \item R2 on $[e_{xy}^{-1}][y^{-1}+e_{xy}]$: trivial prices, gives $[y^{-1}]$.
    \item R2 on $[y^{-1}][y]$: gives $[\emptyset]$.
\end{itemize}

\medskip
\noindent\textbf{Relator 3: $[[x,y],z] \to^* [\emptyset]$ for $\{x,y,z\}\in E_3$, $x<y<z$.}

By the density assumption (Assumption~\ref{assump:dense}), for any
$e \in E$ with $|e| \geq 3$ and any $x,y,z \in e$, the set $\{x,y,z\}$
is itself a hyperedge in $E_3$. It therefore suffices to verify the
reduction for $\{x,y,z\} \in E_3$, which covers all cases.

Here $[[x,y],z]=[e_{xy},z]=e_{xy}^{-1}z^{-1}e_{xy}z$.
Lifting: $[e_{xy}^{-1}][z^{-1}][e_{xy}][z]$, with $z < e_{xy}$.

\begin{itemize}
    \item R1 on $[z^{-1}][e_{xy}]$: the left price collects $[z,e_{xy}]^{-1}=1$ since $\{x,y,z\}\in E_3$ gives $[[x,y],z]=1$ in $G'(\chi)$. Right price $[\emptyset]$. Gives $[z^{-1}+e_{xy}]$.
    \item R2 on $[e_{xy}^{-1}][z^{-1}+e_{xy}]$: trivial prices, gives $[z^{-1}]$.
    \item R2 on $[z^{-1}][z]$: gives $[\emptyset]$.
\end{itemize}

\medskip
\noindent\textbf{Relator 4: $[x,y]\cdot e_{xy}^{-1} \to^* [\emptyset]$.}

The relator in $(A^\pm)^*$ is $x^{-1}y^{-1}xye_{xy}^{-1}$.
Lifting: $[x^{-1}][y^{-1}][x][y][e_{xy}^{-1}]$, with $x < y < e_{xy}$.

\textit{Step 1.} Apply R1 to move $x$ left past $[y^{-1}]$ (addable since $x < y$). Right price $\mathbf{r_x}([y^{-1}])$: the element $y^{-1}$ has $y > x$, contributing $[y,x]^{1\cdot(-1)}=[y,x]^{-1}=e_{xy}$. Left price $[\emptyset]$ (nothing less than $x$ in $[y^{-1}]$). Thus:
$$[y^{-1}][x] \to [x+y^{-1}][e_{xy}][\emptyset] \to^{R3} [x+y^{-1}][e_{xy}]$$

\textit{Step 2.} Apply R2 to $[x^{-1}][x+y^{-1}]$ (with $u=\emptyset$, $x^{-\epsilon}=x^{-1}$, $v=\{y^{-1}\}$). Right price $\mathbf{r_x}([\emptyset])=[\emptyset]$. Left price $\mathbf{l_x}([y^{-1}])=[\emptyset]$ (since $y > x$, no element of $[y^{-1}]$ is less than $x$). After R3:
$$[x^{-1}][x+y^{-1}] \to^* [y^{-1}]$$

Word is now $[y^{-1}][e_{xy}][y][e_{xy}^{-1}]$.

\textit{Step 3.} Apply R1 to move $e_{xy}$ left past $[y^{-1}]$: left price $[y,e_{xy}]^{-1}=1$, right price $[\emptyset]$. Gives $[y^{-1}+e_{xy}]$.

Word is now $[y^{-1}+e_{xy}][y][e_{xy}^{-1}]$.

\textit{Step 4.} Apply R2 to $[y^{-1}+e_{xy}][y]$ (with $u=\{e_{xy}\}$, $x^{-\epsilon}=y^{-1}$, $v=\emptyset$, $x^\epsilon=y$). Right price $\mathbf{r_y}([\{e_{xy}\}])$: $e_{xy}>y$, contributing $[e_{xy},y]^1=[[x,y],y]=1$ (key observation). After R3:
$$[y^{-1}+e_{xy}][y] \to^* [e_{xy}]$$

\textit{Step 5.} Apply R2 to $[e_{xy}][e_{xy}^{-1}]$ (with $u=\emptyset$, $x^{-\epsilon}=e_{xy}$, $x^\epsilon=e_{xy}^{-1}$): trivial prices, gives $[\emptyset]$.

And for the inverse relators:

\medskip
\noindent\textbf{Inverse Relator 1:}
$[[x,y],x]^{-1}=x^{-1}e_{xy}^{-1}xe_{xy}$.
Lifting: $[x^{-1}][e_{xy}^{-1}][x][e_{xy}]$.
\begin{itemize}
  \item R1 on $[x][e_{xy}]$: move $e_{xy}$ into $[x]$. Left price
    $[x,e_{xy}]^{-1}=1$ (key observation). Right price $[\emptyset]$.
    Gives $[x+e_{xy}]$.
  \item R2 on $[e_{xy}^{-1}][x+e_{xy}]$: trivial prices, gives $[x]$.
  \item R2 on $[x^{-1}][x]$: trivial prices, gives $[\emptyset]$.
\end{itemize}

\medskip
\noindent\textbf{Inverse Relator 2:}
$[[x,y],y]^{-1}=y^{-1}e_{xy}^{-1}ye_{xy}$.
Lifting: $[y^{-1}][e_{xy}^{-1}][y][e_{xy}]$.
Identical to Inverse Relator~1 with $y$ in place of $x$.

\medskip
\noindent\textbf{Inverse Relator 3:}
$[[x,y],z]^{-1}=z^{-1}e_{xy}^{-1}ze_{xy}$ for $\{x,y,z\}\in E_3$.
Lifting: $[z^{-1}][e_{xy}^{-1}][z][e_{xy}]$.
Identical structure; left price $[z,e_{xy}]^{-1}=1$ since
$[[x,y],z]=1$ in $G'(\chi)$.

\medskip
\noindent\textbf{Inverse Relator 4:}
$([x,y]e_{xy}^{-1})^{-1}=e_{xy}y^{-1}x^{-1}yx$.
Lifting: $[e_{xy}][y^{-1}][x^{-1}][y][x]$.
\begin{itemize}
  \item R1 on $[x^{-1}][y]$: move $y$ into $[x^{-1}]$. Left price
    $l_y([x^{-1}])=[e_{xy}^{-1}]$ since $x < y$ gives
    $[x,y]^{(-1)(1)}=e_{xy}^{-1}$. Right price $[\emptyset]$.
    Gives $[x^{-1}+y][e_{xy}^{-1}]$.
  \item R2 on $[y^{-1}][x^{-1}+y]$: right price $r_y([y^{-1}])=[\emptyset]$,
    left price $l_y([x^{-1}])=[e_{xy}^{-1}]$. Gives $[e_{xy}^{-1}][x^{-1}]$
    after R3.
  \item Word is now $[e_{xy}][e_{xy}^{-1}][x^{-1}][x]$.
  \item R2 on $[x^{-1}][x]$: trivial prices, gives $[\emptyset]$. Remove
    via R3.
  \item R2 on $[e_{xy}][e_{xy}^{-1}]$: trivial prices, gives
    $[\emptyset]$.
\end{itemize}
\end{proof}

\begin{theorem}[Soundness and Completeness]\label{thm:word_problem}
    Let $w,w'\in \mathcal{A}^*$. Then $\mu(\mathcal{F}(w))=\mu(\mathcal{F}(w'))$ in $G$ if and only if $w$ and $w'$ have the same reduced form.
\end{theorem}

\begin{proof}
$(\Leftarrow)$ Suppose $w$ and $w'$ reduce to the same normal form $\tilde{w}$. Each rewriting step preserves the group element (by Lemma~\ref{lem:interpretation}, each of R1, R2, R3 is a valid group equality). Hence $\mu(\mathcal{F}(w)) = \mu(\mathcal{F}(\tilde{w})) = \mu(\mathcal{F}(w'))$.

$(\Rightarrow)$ Suppose $\mu(\mathcal{F}(w))=\mu(\mathcal{F}(w'))$ in $G'(\chi)$.

By termination (Lemma~\ref{lem:termination}), $w$ and $w'$
reach normal forms $\tilde{w}$ and $\tilde{w}'$ respectively under $\to^*$.
By Lemma~\ref{lem:interpretation}, each rewriting step preserves the group
element, so $\mu(\mathcal{F}(\tilde{w}))=\mu(\mathcal{F}(\tilde{w}'))$.
It remains to show $\tilde{w}=\tilde{w}'$.

Since $\mu(\mathcal{F}(\tilde{w}))=\mu(\mathcal{F}(\tilde{w}'))$, there
exists a finite sequence of elementary moves --- defining relator insertions
and deletions --- connecting $\mathcal{F}(\tilde{w})$ to
$\mathcal{F}(\tilde{w}')$ in $(A^\pm)^*$. Lifting each step via $\sigma$,
this yields a finite sequence of words in $\mathcal{A}^*$:
\[
\tilde{w}=z_0\leftarrow^* p_1\to^* z_1\leftarrow^* p_2\to^*\cdots
\leftarrow^* p_k\to^* z_k=\tilde{w}'
\]
where each $p_i$ is obtained from $z_{i-1}$ by a relator insertion (an
inverse rewriting step, strictly increasing $f$) and reduces to $z_i$
by $\to^*$.

We argue by \textbf{peak reduction} on the maximum value of $f$ along
this sequence. Since $\tilde{w}$ and $\tilde{w}'$ are already in normal
form, any move away from either end strictly increases $f$. Hence there
exists a word $\hat{w}=p_i$ of maximum $f$, with:
\[
z_{i-1}\leftarrow^*\hat{w}\to^* z_i
\]

By local confluence (Lemma~\ref{lem:local_confluence}), there exists a
common descendant $d$ with $z_{i-1}\to^* d$ and $z_i\to^* d$:

\begin{center}
\begin{tikzpicture}[node distance=2cm, auto]
  \node (w) {$\hat{w}$};
  \node (zl) [below left of=w] {$z_{i-1}$};
  \node (zr) [below right of=w] {$z_i$};
  \node (d)  [below right of=zl] {$d$};
  \draw[->] (w)  -- node {$*$}       (zl);
  \draw[->] (w)  -- node[swap] {$*$} (zr);
  \draw[->, dashed] (zl) -- node {$*$}       (d);
  \draw[->, dashed] (zr) -- node[swap] {$*$} (d);
\end{tikzpicture}
\end{center}

Replacing the peak $z_{i-1}\leftarrow^*\hat{w}\to^* z_i$ with the valley
$z_{i-1}\to^* d\leftarrow^* z_i$ yields a new sequence with strictly
smaller maximum $f$. Repeating this peak elimination --- which terminates
since $f$ is well-founded --- produces a sequence with no peaks, i.e.\ a
direct rewriting path $\tilde{w}\to^*\tilde{w}'$. Since $\tilde{w}$ is
already in normal form this path has length zero, giving
$\tilde{w}=\tilde{w}'$ (cf.\ \cite{Sims} \cite{MKS} and \cite{BO93}).

\end{proof}


\begin{lemma}
    Every element of $G$ has a unique reduced representative in $\mathcal{A}^*$, and hence in $(A^{\pm})^*$.
\end{lemma}
\begin{proof}
Existence: every word reduces to some normal form by termination (Lemma~\ref{lem:termination}). Uniqueness: if $w, w'$ are both reduced and $\mu(\mathcal{F}(w)) = \mu(\mathcal{F}(w'))$, then by Theorem~\ref{thm:word_problem} they have the same reduced form, i.e.\ $w = w'$.
\end{proof}

\begin{theorem}[Normal Form Characterization]\label{thm:normal_form}
A word $w = [u_1] \cdots [u_k] \in \mathcal{A}^*$ is in normal form if and only if:
\begin{enumerate}
    \item No $u_i = \emptyset$ \quad (equivalently, R3 does not apply)
    \item For all $i < j$ and all $x^{\epsilon} \in u_j$: $x^{\epsilon}$ is not addable into $u_i$ \quad (equivalently, R1 does not apply)
    \item For all $i$ and all $x^{\epsilon} \in u_i$: $x^{-\epsilon} \notin u_{i+1}$ \quad (equivalently, R2 does not apply)
\end{enumerate}
\end{theorem}

\begin{proof}
($\Rightarrow$) If $w$ is in normal form, no rule applies, so none of the
conditions (1)--(3) can be violated: violation of (1) would trigger R3,
violation of (2) would trigger R1, and violation of (3) would trigger R2.

($\Leftarrow$) Conversely, suppose (1)--(3) all hold. R3 does not apply by
(1). For R1: it requires some $x^\epsilon \in u_j$ to be addable into $u_i$
for $i < j$; this is excluded by (2). For R2: it requires $x^{-\epsilon} \in
u_i$ and $x^\epsilon \in u_{i+1}$ for some $i$ and some $x$; this is excluded
by (3) (taking $j = i+1$ and noting R2 only fires on consecutive pairs with
opposite exponents). Hence no rule applies and $w$ is in normal form.
\end{proof}

\begin{theorem}[Solvability of the Word Problem]\label{thm:solvability}
    The word problem is solvable for partially 2-nilpotent groups.
\end{theorem}

\begin{proof}
Given $w \in (A^{\pm})^*$: (1) lift to $\sigma(w) \in \mathcal{A}^*$; (2) apply rewriting rules until none applies, obtaining normal form $\tilde{w}$ (terminates by Lemma~\ref{lem:termination}); (3) $w$ represents the identity if and only if $\tilde{w} = [\emptyset]$, by Theorem~\ref{thm:word_problem}.
\end{proof}

\end{section}

\printbibliography

\end{document}